\section{Protocol Robustness and Integrity}
\label{sec:robust}

\begin{guidance}[Protocol robustness and integrity (10 marks)]
Show \textbf{framing}, \textbf{specific NACKs for malformed input}, and
\textbf{graceful recovery} from partial or garbage data. A table that maps each
kind of bad input to your mechanism, its response and the test that checks it
works well. Also explain why the server can never block forever (bounded
loops, timeouts, no blocking waits in the main loop). Give measured numbers
where you have them, e.g. the timeout you observed.
\end{guidance}

The serial game server adopts a defense-in-depth model that combines deterministic interrupt-driven frame assembly, non-blocking inter-byte silence timeouts, and multi-tier semantic gating to guarantee fail-safe operation against corrupt, truncated, or malicious host inputs.

\begin{table*}[t]
\centering
\caption{Fault taxonomy, detection and mitigation mechanisms, server responses, and automated validation tests.}
\label{tab:threats}
\small
\begin{tabular}{@{}L{2.5cm}L{6.2cm}L{2.5cm}L{5.6cm}@{}}
\toprule
\textbf{Input Fault} & \textbf{Detection and Handling Mechanism} & \textbf{Server Response} & \textbf{Automated Test Identifier} \\
\midrule
Partial / Truncated Frame & Inter-byte silence timer ($\Delta t > 500$\,ms) in main loop resets \texttt{rx\_index} & Silent discard; buffer cleared & \texttt{test\_incomplete\_frame\_}\allowbreak\texttt{timeout\_recovery} \\
\addlinespace[2pt]
Oversized / Flooded Frame & Boundary guard: \texttt{rx\_index < MAX\_FRAME\_LEN - 1} resets index on saturation & Silent discard; buffer reset & \texttt{test\_buffer\_overflow\_}\allowbreak\texttt{protection} \\
\addlinespace[2pt]
Garbage Bytes & Header validation failure in \texttt{Proto\_ParseFrame()} & \texttt{NACK: BADCMD} & \texttt{test\_malformed\_command\_}\allowbreak\texttt{rejection} \\
\addlinespace[2pt]
Malformed Syntax & Delimiter, coordinate bounds, or player ID check failure & \texttt{NACK: BADCMD} & \texttt{test\_malformed\_shoot\_}\allowbreak\texttt{commands} \\
\addlinespace[2pt]
Out-of-Turn Move & Active player mismatch (\texttt{cmd->player != current\_turn}) & \texttt{NACK: NOTURN} & \texttt{test\_player\_two\_cannot\_}\allowbreak\texttt{shoot\_first} \\
\addlinespace[2pt]
Out-of-State Move & State gate in \texttt{Game\_ProcessCommand()} dispatch & \texttt{NACK: NOGAME} & \texttt{test\_place\_command\_}\allowbreak\texttt{rejected\_during\_gameplay} \\
\addlinespace[2pt]
Illegal Fleet Placement & Collision, diagonal, or size violation in atomic placement & \texttt{NACK: OCCUPIED} / \texttt{NACK: BADCOORDS} & \texttt{test\_overlapping\_ship\_}\allowbreak\texttt{rejected} \\
\addlinespace[2pt]
Repeated Target & Bitwise hit mask verification ($\text{cell} \ \& \ 0\text{x}80 \ne 0$) & \texttt{NACK: REPEAT} & \texttt{test\_repeat\_shot\_}\allowbreak\texttt{rejected\_and\_retains\_turn} \\
\bottomrule
\end{tabular}
\end{table*}

\subsection{Framing, Delimiter Parsing, and Specific Rejections}
Communication over the asynchronous UART channel is packetized into ASCII line frames terminated by a newline character (\texttt{\textbackslash n}), with carriage returns (\texttt{\textbackslash r}) stripped on ingress by \texttt{Proto\_ProcessByte()} to preserve cross-platform compatibility across Windows CRLF and POSIX LF conventions. 

When a newline is received, the interrupt sets \texttt{frame\_ready = 1}, alerting the foreground superloop to parse the null-terminated payload via \texttt{Proto\_ParseFrame()}. Every incoming message is strictly classified against expected grammatical structures:
\begin{enumerate}
    \item \textbf{Token Prefix Verification}: The command verb must match known tokens (\texttt{CONNECT}, \texttt{DISCONNECT}, \texttt{NEW: }, \texttt{STATE: }, \texttt{FLEET: }, \texttt{SHOOT: }). Any unrecognized string or case alteration (e.g.\ \texttt{shoot:}) yields an immediate \texttt{NACK: BADCMD}.
    \item \textbf{Delimiter and Bounds Enforcement}: For parameterized frames, comma placement, string lengths, and alphanumeric boundaries are checked. For example, \texttt{SHOOT: 1,A1} requires length $\ge 10$, player $\in \{1, 2\}$, column $\in [\text{`A'}, \text{`J'}]$, and row $\in [1, 10]$. Omission of spaces (e.g.\ \texttt{SHOOT:1,A1}), missing commas (\texttt{SHOOT: 1A1}), or out-of-grid coordinates (\texttt{SHOOT: 1,Z99}) immediately fail validation, triggering \texttt{NACK: BADCMD}.
    \item \textbf{Semantic State and Rule Rejections}: Frames with valid syntax but invalid context are rejected with contextual error codes (\texttt{NACK: NOGAME}, \texttt{NACK: NOTURN}, \texttt{NACK: OCCUPIED}, or \texttt{NACK: REPEAT}), leaving internal state invariants intact as proven in Section~\ref{sec:authority}.
\end{enumerate}

\subsection{Partial Frames, Inter-Byte Silence, and Overflow Recovery}
A robust serial interface must tolerate physical cable disconnections, host process termination, and electrical noise bursts without latching into an unrecoverable framing state.

\subsubsection{Inter-Byte Silence Timeout}
If a client transmits a partial frame (e.g.\ \texttt{CON} or \texttt{SHOOT: 1,}) and ceases transmission, the partial bytes remain buffered in \texttt{rx\_buffer}. Without an eviction policy, subsequent valid commands would append to the fragment (forming \texttt{CONCONNECT\textbackslash n}), causing perpetual framing errors. 

To resolve this, the server implements an inter-byte silence timeout. Rather than measuring total frame transmission time (which would penalize long multi-coordinate placement frames), the foreground loop tracks the elapsed silence since the last received byte:
\begin{equation}
\Delta t_{\text{silence}} = t_{\text{now}} - t_{\text{last\_rx}}
\end{equation}
Whenever $0 < \texttt{rx\_index} < \text{MAX\_FRAME\_LEN}$ and $\Delta t_{\text{silence}} > \text{FRAME\_TIMEOUT}$ (configured to 500\,ms in \texttt{main.h}), the buffer is automatically flushed:
\begin{lstlisting}[language=C]
if (rx_index == 0) {
  last_rx_time = now;
} else if (now - last_rx_time > FRAME_TIMEOUT) {
  rx_index = 0;
  last_rx_time = now; // Reset timer after clearing
}
\end{lstlisting}

At the nominal baud rate of 115,200\,bps (8-N-1 framing), one byte spans:
\begin{equation}
T_{\text{byte}} = \frac{10\,\text{bits}}{115{,}200\,\text{baud}} \approx 86.81\,\mu\text{s}
\end{equation}
A full 64-byte frame is transmitted in approximately 5.56\,ms. The selected 500\,ms threshold corresponds to $\approx 5{,}760$ byte transmission intervals. This provides a generous safety margin that easily absorbs host operating system scheduling delays and USB-to-UART CDC buffer latency (typically $\le 16$\,ms), while ensuring rapid recovery within half a second of a transmission fault. Automated verification in \texttt{test\_incomplete\_frame\_timeout\_recovery} injects \texttt{"CON"}, sleeps for 600\,ms, and confirms that a subsequent \texttt{"CONNECT\textbackslash n"} succeeds with \texttt{"OK"}.

\subsubsection{Buffer Overflow Protection and Resynchronization}
To protect against denial-of-service conditions or unbounded memory corruption from malformed high-throughput serial floods, \texttt{Proto\_ProcessByte()} enforces a strict boundary guard against \texttt{MAX\_FRAME\_LEN} (64 bytes). If an incoming stream delivers 63 consecutive non-newline bytes, the buffer pointer is automatically reset to zero:
\begin{lstlisting}[language=C]
if (rx_index < (MAX_FRAME_LEN - 1)) {
    rx_buffer[rx_index++] = (char)byte;
} else {
    // Saturated frame: clear buffer to recover
    rx_index = 0;
}
\end{lstlisting}
This prevents index overruns into adjacent SRAM variables. Furthermore, the parser naturally resynchronizes on the next newline character (\texttt{\textbackslash n}): the discarded run is nullified, and the next incoming byte is treated as the start of a clean frame. Automated test \texttt{test\_buffer\_overflow\_protection} exercises this safeguard by streaming 160 continuous bytes of garbage (\texttt{"GARBAGE\_" * 20 + "\textbackslash n"}); the firmware absorbs the flood without a HardFault, immediately processing subsequent valid transactions.

\subsection{Non-Blocking Architecture and Hang Prevention}
Embedded reliability dictates that the microcontroller must never lock up or hang indefinitely, regardless of client malfunction or electrical noise. The server achieves unconditional forward progress through three architectural guarantees:
\begin{enumerate}
    \item \textbf{Absence of Unbounded Loops}: Every loop construct across the firmware operates under statically determinable, compile-time upper bounds. Grid zeroing and fleet placements iterate across a maximum of 100 cells ($10 \times 10$) and 5 ships of length $\le 5$ (at most 25 cells checked). The probabilistic AI search scans the grid in a fixed $10 \times 10$ loop with zero runtime dynamic branch allocation.
    \item \textbf{Non-Blocking Superloop Execution}: The primary execution cycle in \texttt{main.c} avoids all blocking polling loops. Frame parsing, state updates, LED rendering, and buzzer cadence generation execute asynchronously based on non-blocking millisecond tick evaluations (\texttt{HAL\_GetTick()}).
    \item \textbf{Interrupt-Driven Reception and Monitored Transmissions}: Serial ingestion occurs entirely within the non-blocking \texttt{USART1} RX interrupt. Outbound responses use bounded hardware polling (\texttt{HAL\_UART\_Transmit()}), where string lengths are restricted to small fixed buffers ($\le 64$ bytes), completing transmission in under 6\,ms.
\end{enumerate}
Consequently, the worst-case foreground iteration latency is bounded under 12\,ms, guaranteeing perpetual responsiveness to host commands and physical inputs.
